Offline-to-online reinforcement learning can reuse logged experience before interacting with the environment, but constraints that help under logged replay may not transfer cleanly to online adaptation. We study this constraint-transfer gap in a focused D4RL MuJoCo replay setting. Plain TD3+BC is weak on \texttt{hopper-medium-replay-v2}: it reaches 22.43 normalized score after 50k offline steps, while ReBRAC-lite reaches 34.48 and SSAR reaches 38.56 final / 43.97 best under the same 50k screen. Mechanism experiments suggest that SSAR's IQL-qv trusted-action selection is a key driver: removing trusted selection collapses performance, while an aligned teacher-label selector, ATLAS, reaches 69.97 and 68.11 final score on hopper seed0/seed1 at 100k steps. Shuffling labels collapses ATLAS to 18.78, indicating that label alignment matters rather than arbitrary weighting. However, the eval20 offline-to-online panel reveals the limitation: TD3+BC improves online, while stronger teacher labels often lose final score when carried forward naively. A Q-filtered ATLAS gate fixes this on seed0 (48.41 final), but seed1 shows large Hopper O2O variance and no stable dominance over TD3+BC. Our evidence supports a narrow workshop-style finding: trusted-action labels are useful offline supervision under low-quality replay, but online fine-tuning needs adaptive constraint transfer rather than fixed teacher regularization.
